% GENERATED FILE. Edit canonical lessons, then run npm run generate:books.
% canonical-source-hash: fnv1a64:cf4b90fa94eb6e42
% canonical-lessons: GU-C467-vartan, GU-C467-adab, GU-C467-nakhra, GU-C467-caritrya, GU-C467-lajvu

\chapter{Good Manners}
\label{ch:gu-good-manners}
\hlchaptermodality{467}

\begin{chapteropening}
\textbf{By the end of this chapter:} \emph{I can talk about good and bad behaviour.}

Complete the last lesson of chapter 467: I can talk about good and bad behaviour.
\end{chapteropening}

\section[\texorpdfstring{\gu{વર્તન} (vartan)}{વર્તન (vartan)}]{\gu{વર્તન} (vartan) — behaviour}

\label{lesson:GU-C467-vartan}



\begin{quote}
Before the new one: say the Gujarati for grit, then the Gujarati for to wobble.
\end{quote}

\subsection*{You'll want to know: \gu{વર્તન}}
\textbf{\gu{વર્તન}} — \emph{vartan} — \textquotedblleft{}behaviour\textquotedblright{}.

\begin{quote}
\gu{મહેમાનો} \gu{સાથે} \gu{તમારું} \gu{વર્તન} \gu{સારું} \gu{હતું}.
\end{quote}

\begin{quote}
\emph{mahemāno sāthe tamārũ vartan sārũ hatũ.}
\end{quote}

\begin{quote}
You behaved well with the guests.
\end{quote}

A neuter noun.

\begin{cousinweb}
\textbf{\gu{વર્તન}} (\emph{vartan}) is a tatsama, Sanskrit \emph{vartana}, which Monier-Williams gives as \textquotedblleft{}abiding\textquotedblright{} and \textquotedblleft{}setting in motion\textquotedblright{}; Gujarati uses it for the way a person conducts himself.
\end{cousinweb}

\begin{grammarlens}[title={What you've built}]
One more word for describing what you lived and felt.
\end{grammarlens}

\subsection*{Your turn}
\begin{itemize}
\raggedright
  \item \emph{Say it:} \emph{vartan}
  \item \emph{Say it:} \emph{vartan}, once more
  \item \emph{Say it:} \emph{mahemāno sāthe tamārũ vartan sārũ hatũ.} — You behaved well with the guests.
  \item \emph{Recall:} say the Gujarati for grit, then the Gujarati for to wobble, then say \emph{vartan} again
\end{itemize}

\begin{tcolorbox}[breakable,colback=teal!4,colframe=teal!35!black,title={Before you move on}]
What does \gu{વર્તન} mean? (\textquotedblleft{}Behaviour\textquotedblright{}.) Say it once more.
\end{tcolorbox}
\section[\texorpdfstring{\gu{અદબ} (adab)}{અદબ (adab)}]{\gu{અદબ} (adab) — respectful manners}

\label{lesson:GU-C467-adab}



\begin{quote}
Before the new one: say the Gujarati for to wobble, then the Gujarati for behaviour.
\end{quote}

\subsection*{You'll want to know: \gu{અદબ}}
\textbf{\gu{અદબ}} — \emph{adab} — \textquotedblleft{}respectful manners\textquotedblright{}.

\begin{quote}
\gu{બાળકોએ} \gu{દાદા} \gu{સાથે} \gu{અદબથી} \gu{બોલવું} \gu{જોઈએ}.
\end{quote}

\begin{quote}
\emph{bāḷkoe dādā sāthe adabthī bolvũ joīe.}
\end{quote}

\begin{quote}
Children should speak to their grandfather respectfully.
\end{quote}

A feminine noun.

\begin{cousinweb}
\textbf{\gu{અદબ}} (\emph{adab}) is usually traced to Arabic \emph{adab}, \textquotedblleft{}good manners, culture\textquotedblright{}, which reached Gujarati through Persian.
\end{cousinweb}

\begin{grammarlens}[title={What you've built}]
One more word for describing what you lived and felt.
\end{grammarlens}

\subsection*{Your turn}
\begin{itemize}
\raggedright
  \item \emph{Say it:} \emph{adab}
  \item \emph{Say it:} \emph{adab}, once more
  \item \emph{Say it:} \emph{bāḷkoe dādā sāthe adabthī bolvũ joīe.} — Children should speak to their grandfather respectfully.
  \item \emph{Recall:} say the Gujarati for to wobble, then the Gujarati for behaviour, then say \emph{adab} again
\end{itemize}

\begin{tcolorbox}[breakable,colback=teal!4,colframe=teal!35!black,title={Before you move on}]
What does \gu{અદબ} mean? (\textquotedblleft{}Respectful manners\textquotedblright{}.) Say it once more.
\end{tcolorbox}
\section[\texorpdfstring{\gu{નખરાં} (nakhrā̃)}{નખરાં (nakhrā̃)}]{\gu{નખરાં} (nakhrā̃) — airs; fussy demands}

\label{lesson:GU-C467-nakhra}



\begin{quote}
Before the new one: say the Gujarati for behaviour, then the Gujarati for respectful manners.
\end{quote}

\subsection*{You'll want to know: \gu{નખરાં}}
\textbf{\gu{નખરાં}} — \emph{nakhrā̃} — \textquotedblleft{}airs; fussy demands\textquotedblright{}.

\begin{quote}
\gu{ભોજન} \gu{વખતે} \gu{નાની} \gu{બહેન} \gu{ઘણાં} \gu{નખરાં} \gu{કરે} \gu{છે}.
\end{quote}

\begin{quote}
\emph{bhojan vakhte nānī bahen ghaṇā̃ nakhrā̃ kare che.}
\end{quote}

\begin{quote}
At mealtimes my little sister makes a lot of fuss.
\end{quote}

A neuter noun, used in the plural.

\begin{cousinweb}
\textbf{\gu{નખરાં}} (\emph{nakhrā̃}) is usually traced to Persian \emph{naxra}, \textquotedblleft{}coquetry, airs\textquotedblright{}.
\end{cousinweb}

\begin{grammarlens}[title={What you've built}]
One more word for describing what you lived and felt.
\end{grammarlens}

\subsection*{Your turn}
\begin{itemize}
\raggedright
  \item \emph{Say it:} \emph{nakhrā̃}
  \item \emph{Say it:} \emph{nakhrā̃}, once more
  \item \emph{Say it:} \emph{bhojan vakhte nānī bahen ghaṇā̃ nakhrā̃ kare che.} — At mealtimes my little sister makes a lot of fuss.
  \item \emph{Recall:} say the Gujarati for behaviour, then the Gujarati for respectful manners, then say \emph{nakhrā̃} again
\end{itemize}

\begin{tcolorbox}[breakable,colback=teal!4,colframe=teal!35!black,title={Before you move on}]
What does \gu{નખરાં} mean? (\textquotedblleft{}Airs; fussy demands\textquotedblright{}.) Say it once more.
\end{tcolorbox}
\section[\texorpdfstring{\gu{ચારિત્ર્ય} (cāritrya)}{ચારિત્ર્ય (cāritrya)}]{\gu{ચારિત્ર્ય} (cāritrya) — moral character}

\label{lesson:GU-C467-caritrya}



\begin{quote}
Before the new one: say the Gujarati for respectful manners, then the Gujarati for airs.
\end{quote}

\subsection*{You'll want to know: \gu{ચારિત્ર્ય}}
\textbf{\gu{ચારિત્ર્ય}} — \emph{cāritrya} — \textquotedblleft{}moral character\textquotedblright{}.

\begin{quote}
\gu{સારું} \gu{ચારિત્ર્ય} \gu{પૈસા} \gu{કરતાં} \gu{મોટું} \gu{છે}.
\end{quote}

\begin{quote}
\emph{sārũ cāritrya paisā kartā̃ moṭũ che.}
\end{quote}

\begin{quote}
Good character is worth more than money.
\end{quote}

A neuter noun.

\begin{cousinweb}
\textbf{\gu{ચારિત્ર્ય}} (\emph{cāritrya}) is a tatsama, Sanskrit \emph{cāritrya}, \textquotedblleft{}good conduct\textquotedblright{}, as Monier-Williams records it.
\end{cousinweb}

\begin{grammarlens}[title={What you've built}]
One more word for describing what you lived and felt.
\end{grammarlens}

\subsection*{Your turn}
\begin{itemize}
\raggedright
  \item \emph{Say it:} \emph{cāritrya}
  \item \emph{Say it:} \emph{cāritrya}, once more
  \item \emph{Say it:} \emph{sārũ cāritrya paisā kartā̃ moṭũ che.} — Good character is worth more than money.
  \item \emph{Recall:} say the Gujarati for respectful manners, then the Gujarati for airs, then say \emph{cāritrya} again
\end{itemize}

\begin{tcolorbox}[breakable,colback=teal!4,colframe=teal!35!black,title={Before you move on}]
What does \gu{ચારિત્ર્ય} mean? (\textquotedblleft{}Moral character\textquotedblright{}.) Say it once more.
\end{tcolorbox}
\section[\texorpdfstring{\gu{લાજવું} (lājvũ)}{લાજવું (lājvũ)}]{\gu{લાજવું} (lājvũ) — to be ashamed}

\label{lesson:GU-C467-lajvu}



\begin{quote}
Before the new one: say the Gujarati for airs, then the Gujarati for moral character.
\end{quote}

\subsection*{You'll want to know: \gu{લાજવું}}
\textbf{\gu{લાજવું}} — \emph{lājvũ} — \textquotedblleft{}to be ashamed\textquotedblright{}.

\begin{quote}
\gu{ખોટું} \gu{બોલ્યા} \gu{પછી} \gu{તે} \gu{લાજી} \gu{ગયો}.
\end{quote}

\begin{quote}
\emph{khoṭũ bolyā pachī te lājī gayo.}
\end{quote}

\begin{quote}
After telling a lie he felt ashamed.
\end{quote}

Past (she …): \gu{તે} \gu{લાજી} \gu{ગઈ}, \emph{te lājī gaī}.

\begin{cousinweb}
\textbf{\gu{લાજવું}} (\emph{lājvũ}) is inherited from Sanskrit \emph{lajjatē}, \textquotedblleft{}is ashamed\textquotedblright{}; CDIAL lists the Gujarati verb (CDIAL 10909).
\end{cousinweb}

\begin{grammarlens}[title={What you've built}]
One more word for describing what you lived and felt.
\end{grammarlens}

\subsection*{Your turn}
\begin{itemize}
\raggedright
  \item \emph{Say it:} \emph{lājvũ}
  \item \emph{Say it:} \emph{lājvũ}, once more
  \item \emph{Say it:} \emph{khoṭũ bolyā pachī te lājī gayo.} — After telling a lie he felt ashamed.
  \item \emph{Recall:} say the Gujarati for airs, then the Gujarati for moral character, then say \emph{lājvũ} again
\end{itemize}

\begin{tcolorbox}[breakable,colback=teal!4,colframe=teal!35!black,title={Before you move on}]
What does \gu{લાજવું} mean? (\textquotedblleft{}To be ashamed\textquotedblright{}.) Say it once more.
\end{tcolorbox}
